\section{Conclusion}
\label{sec:conclusion}

There were many stages throughout the creation of this project. Namely, an analysis of the original dataset, designing the data warehouse, formulating an ETL script, formulating queries, optimizing queries, and finally making a web app to add a user interface for everything. And as a natural result of creating a project, these concepts and processes all became more understandable after the project.

Switching from an OLTP thought process to OLAP was one of the more difficult parts in this project to adapt to, and was especially seen when formulating the data warehouse. In OLTP, and in CCINFOM, we worked primarily with normalized schemas and simple operations that fetch data as it is written. However, in building the data warehouse and ETL for the OLAP application, it's necessary to keep in consideration that this is ultimately not meant to be a normalized schema, but rather a schema that is meant to aid in analysis. Denormalization and derived columns are methods that make queries faster and simpler. In schema designing, it is necessary to shift mindset from making the most organized schema, but rather one that makes analyzation more efficient.

In creating this project, the necessity of query optimization is once again highlighted. Even though the dataset given is not necessarily massive, it is large enough that there were noticeable wait times between queries when first writing them. Common Table Expressions (CTEs), materialization, and indexing were used to help in optimization. Besides increasing the speed of the queries, optimization also encourages thinking about queries in a more analytical manner. It encourages spotting redundancies and bottlenecks, and once removed, the query is often more readable as an additional bonus.

There is no shortage of data in the world, but a great majority of it goes unused. This Dutch railway application was build with the goal of providing insights that help people find underlying problems in the Dutch railway that can be amended. In learning how to build better data warehouses, ETL scripts, and queries, we become better equipped to help make our abundance of data into information that is valuable in the long run.
